\section{Optimization Problem}
\subsection{Formulation of the optimization problem}\;\\
Using the vehicle model introduced in the previous section, the minimum-lap-time problem can be formulated as a continuous-time optimal control problem. The decision variables are the state trajectory $z(\cdot)$, the control trajectory $u(\cdot)$, and the final time $T$.

The control vector
\begin{align*}
    u(t) = \begin{bmatrix} a(t) \\ \delta(t) \end{bmatrix}
\end{align*}
contains the longitudinal acceleration and the steering angle.

The objective is to minimize the total time required to complete the track. 
The resulting optimal control problem is
\begin{align}
    \underset{z(\cdot),\,u(\cdot),\,T}{\operatorname{minimize}}
    \quad & T
    \label{eq:objective}
    \\[0.5em]
    \text{subject to}
    \quad
    & \dot{z}(t)
    =
    f\bigl(z(t),u(t)\bigr),
    && t\in[0,T],
    \label{eq:dynamics}
    \\[0.3em]
    & -\frac{w\bigl(s(t)\bigr)-b}{2}
    \leq e_y(t)
    \leq
    \frac{w\bigl(s(t)\bigr)-b}{2},
    && t\in[0,T],
    \label{eq:track_constraint}
    \\[0.3em]
    & -a_{\min}
    \leq a(t)
    \leq a_{\max},
    && t\in[0,T],
    \label{eq:acceleration_constraint}
    \\[0.3em]
    & -\delta_{\max}
    \leq \delta(t)
    \leq \delta_{\max},
    && t\in[0,T],
    \label{eq:steering_constraint}
    \\[0.3em]
    & 0
    \leq v(t)
    \leq v_{\max},
    && t\in[0,T],
    \label{eq:velocity_constraint}
    \\[0.3em]
    & -a_{y,\max}
    \leq
    \frac{v(t)^2}{L}\tan\bigl(\delta(t)\bigr)
    \leq
    a_{y,\max},
    && t\in[0,T],
    \label{eq:lateral_acceleration_constraint}
    \\[0.3em]
    & 1-\kappa\bigl(s(t)\bigr)e_y(t)>0,
    && t\in[0,T],
    \label{eq:frenet_validity_constraint}
    \\[0.3em]
    & s(0)=0,
    \qquad
    s(T)=S,
    \label{eq:position_boundary_conditions}
    \\[0.3em]
    & v(0)=0,
    \qquad
    e_y(0)=0,
    \qquad
    e_\psi(0)=0,
    \label{eq:initial_conditions}
    \\[0.3em]
    & T>0.
    \label{eq:time_constraint}
\end{align}

The parameters appearing in the optimal control problem are defined as follows:
\begin{description}
    \item[\(w(s)\)] track width at centerline position \(s\),
    \item[\(b\)] vehicle width,
    \item[\(L\)] vehicle wheelbase,
    \item[\(a_{\min}>0\)] magnitude of the maximum permissible deceleration,
    \item[\(a_{\max}>0\)] maximum permissible acceleration,
    \item[\(\delta_{\max}\)] maximum absolute steering angle,
    \item[\(v_{\max}\)] maximum permissible velocity,
    \item[\(a_{y,\max}\)] maximum permissible lateral acceleration,
    \item[\(S\)] total centerline length of the track.
\end{description}

The objective function~\eqref{eq:objective} minimizes the final time $T$.
The dynamic constraint~\eqref{eq:dynamics} ensures that the state trajectory
is consistent with the kinematic bicycle model.

The track constraint~\eqref{eq:track_constraint} restricts the center of the
vehicle to the usable track width. The vehicle width $b$ is subtracted from
the track width so that the complete vehicle, rather than only its reference
point, remains inside the track boundaries.

The acceleration, steering, and velocity constraints
\eqref{eq:acceleration_constraint}--\eqref{eq:velocity_constraint} represent
the longitudinal and kinematic limitations of the vehicle. The lateral
acceleration constraint~\eqref{eq:lateral_acceleration_constraint} limits the
cornering acceleration generated by the steering input.

The boundary conditions require the vehicle to start at the beginning of the
track and reach the total centerline distance $S$. In the present formulation,
the vehicle starts from rest at the centerline with zero heading error.


\subsection{Extension: Static obstacle avoidance.}\;\\
So far, the optimization problem only considers the track boundaries and vehicle limitations. In practical racing scenarios, however, additional objects may occupy parts of the track. To account for such situations, the formulation can be extended by introducing static obstacles represented in Frenet coordinates.

The positions of the obstacles are assumed to be known in advance and are specified in Frenet coordinates. For each obstacle $i=1,\ldots,N_{\mathrm{obs}}$, the centerline position is denoted by $s_i$ and the lateral offset by $e_{y,i}$. Furthermore, each obstacle is assigned a length $l_i$ and a width $w_i$.

To account for the finite dimensions of both the obstacle and the vehicle, together with an additional safety margin, each obstacle is surrounded by an elliptical safety region. The corresponding semi-axes in the longitudinal and lateral directions are defined as
\begin{align*}
    & r_{s,i}=\frac{l_i}{2}+\frac{L}{2}+d_{\mathrm{safe}},
    \\
    & r_{e_y,i}=\frac{w_i}{2}+\frac{b}{2}+d_{\mathrm{safe}},
\end{align*}
where $L$ denotes the vehicle wheelbase, $b$ its width, and $d_{\mathrm{safe}}$ the prescribed safety margin.

For every obstacle, the vehicle is required to remain outside the corresponding safety region. This requirement is enforced by introducing the nonlinear constraint
\begin{align}
\left(\frac{s-s_i}{r_{s,i}}\right)^2
+
\left(\frac{e_y-e_{y,i}}{r_{e_y,i}}\right)^2
\ge 1,
\qquad i=1,\ldots,N_{\mathrm{obs}}.
\end{align}

Since the optimization problem is solved on a discrete time grid, obstacle avoidance constraints are additionally evaluated at the midpoint of every integration interval. This improves collision detection between consecutive discretization nodes while avoiding the computational expense of using a much finer discretization.

This extension illustrates how additional practical constraints can be integrated into the optimal control formulation without modifying the underlying vehicle model.